\section{补充例子与反例}
\label{app:examples-counterexamples}

\subsection{循环链否定普遍 LC}

\begin{counterexample}[有限性不推出低复杂度]
\label{counter:cycle}
令 $B_N=\mathbb Z/N\mathbb Z$，确定性 kernel 为
$K_N(x,x+1)=1$，任务为 $g_N(x)=\one_{\{x=0\}}$。若 $\tau<1/2$，任何任务保真且闭合缺陷不超过 $\tau$ 的分区都只能是离散分区。因此
\[
c_{K_N}^{\tau}(g_N)=N,
\]
系统族不满足 LC。
\end{counterexample}

\begin{proof}
若 $s(x)=s(y)$，两条投影后的一步分布 $\delta_{s(x+1)}$ 与 $\delta_{s(y+1)}$ 都在 TV 距离 $\tau$ 内接近同一宏观行。因此它们之间距离不超过 $2\tau<1$。两个 Dirac 分布的 TV 距离只能是 $0$ 或 $1$，故 $s(x+1)=s(y+1)$；迭代得到 $s(x+t)=s(y+t)$ 对所有 $t$ 成立。若 $x\ne y$，可选 $t$ 使其中一个状态等于 $0$、另一个不等于 $0$，与任务保真矛盾。
\end{proof}

\subsection{闭合不需要非平凡状态对称}

\begin{example}[无交换对称的闭合商]
取 $B=\{1,2,3\}$，分区 $\pi=\{\{1,2\},\{3\}\}$，并令
\[
K=\begin{pmatrix}
0.1&0.6&0.3\\
0.2&0.5&0.3\\
0.1&0.2&0.7
\end{pmatrix}.
\]
前两行到两个分区块的概率都为 $(0.7,0.3)$，所以 $\pi$ 精确闭合；但交换状态 $1,2$ 并不保持完整 kernel。由此可见，微观群等变性是方便的充分证书，不是闭合的必要条件。
\end{example}

\subsection{逐环境最优与共享最优}

\begin{counterexample}[逐环境规律不等于共享规律]
令 $B=\{0,1\}\times\{a,b\}$，$Y=\{0,1\}$，$q(y,z)=y$，任务取 $g=q$。构造两个 kernels：$K^{(0)}$ 的下一宏观状态恒为 $0$，$K^{(1)}$ 的下一宏观状态恒为 $1$；纤维内的第二坐标可任意按固定规则更新。则 $q$ 对两者都精确闭合，分别对应
$L^{(0)}(y,\cdot)=\delta_0$ 与 $L^{(1)}(y,\cdot)=\delta_1$，故
\[
\sup_{i\in\{0,1\}}\min_L
d_\infty(K^{(i)}Q_q,Q_qL)=0.
\]
但任意共享 $L$ 的每一行都必须同时逼近 $\delta_0$ 与 $\delta_1$，由三角不等式其最坏误差至少 $1/2$，而均匀分布达到该值。因此
\[
\min_L\sup_{i\in\{0,1\}}
d_\infty(K^{(i)}Q_q,Q_qL)=\frac12.
\]
\end{counterexample}

\subsection{破缺相不推出轨道压缩}

\begin{counterexample}[相形成与低复杂度分离]
令 $B_N=\{+,-\}\times\{1,\ldots,N\}$，$G_N=\mathbb Z_2$ 交换第一坐标，$K_N=\Id_{B_N}$，$\mu_N=\delta_{(+,1)}$，并取残余子群 $H_N=\{e\}$。则 $K_N$ 对 $G_N$ 精确等变，$\mu_N$ 精确平稳，且相对 $G_N$ 存在精确分离的破缺相；但是
\[
\frac{|B_N/H_N|}{|B_N|}=1.
\]
因此精确平稳、精确大群等变和非平凡破缺相同时成立，也不推出轨道商具有低相对复杂度。
\end{counterexample}

这些反例分别排除四种过强解读：有限性不保证 LC；对称不是闭合的必要条件；逐环境闭合不等于共享宏观核；相形成不等于宏观压缩。
